\documentclass[UTF8]{ctexart}
\usepackage[a4paper,margin=2.6cm]{geometry}
\begin{document}

\section*{摘要}
这篇论文讨论的是：为什么 Dijkstra 算法不仅是最经典的单源最短路算法之一，而且在一种更细的、超出最坏情况的分析框架下，也可以被看成是非常自然、非常接近最优的算法。

作者把重点放在“距离顺序”这件事上。Dijkstra 的过程本质上是在不断确定顶点到源点的距离顺序，而这类问题的复杂度不只取决于图有多大，也取决于输入本身的结构。论文提出一种比传统最坏情况分析更强的视角，用来说明 Dijkstra 的性能在很多输入上都具有普遍的最优性。

为了让这种分析真正对应到算法实现，作者还构造了一个适合 Dijkstra 的堆结构。这个堆不是只追求抽象上界，而是要能支持大量 \texttt{decrease-key} 操作，并且同时保留类似 working-set 的行为。论文的技术价值主要就在于：它把这种“普遍最优”的分析和一个可实现的优先队列结合了起来。

\end{document}
